\section{Introduction}

Payment fraud is a well-suited application of supervised learning. Labeled transactions are plentiful, the cost of an error is measurable, and a model's score can be converted directly into a decision to block or review a transaction. Two public datasets in particular, the credit card dataset from Universit\'e Libre de Bruxelles (ULB)~\cite{dalpozzolo2015calibrating} and the IEEE-CIS competition dataset~\cite{ieeecis2019fraud}, now underpin a large body of published comparisons.

A deployed fraud model is trained on past transactions and scores future ones. Much of the published work evaluates models differently. Transactions are shuffled and split at random into training and test sets, or into cross-validation folds~\cite{esenogho2022neural,alfaiz2022enhanced,mim2024soft,du2024novel}, so that the model is tested on transactions that are interleaved in time with the ones it was trained on. Some published pipelines describe rebalancing the classes with SMOTE~\cite{chawla2002smote} or similar resampling before partitioning the data~\cite{ileberi2021performance,khalid2024enhancing}, which places synthetic neighbors of test fraud cases in the training set. These are two forms of data leakage~\cite{kaufman2012leakage,kapoor2023leakage}, and the accuracies above 99\% that often result say little about how a model would perform in production.

Previous work has noted that random splits produce optimistic estimates for fraud models~\cite{grover2022fraud,breskuviene2024enhancing,hayat2025data}, but a controlled measurement of the size of the effect, its dependence on the model, and its cause has been lacking. In this paper we train five standard classifiers on both public datasets under a random split and under several chronological protocols, holding the tuning budget, the training-set size and all other settings fixed, so that the difference in scores can be attributed to the split.

The contributions of the paper are as follows.
\begin{itemize}[leftmargin=1.5em,itemsep=2pt]
\item \textbf{A like-for-like measurement of split inflation.} On IEEE-CIS, moving from a random to a chronological split lowers \prauc{} by \num{ieee.min.rel.pr}--\num{ieee.max.rel.pr}\% for every model, and the effect is consistent across rolling windows. We add two controls that score both protocols on the same transactions. They provide strong evidence that temporal interleaving, rather than a harder test period, explains most of the IEEE-CIS gap, and they show that most of the apparent gap on ULB is due to a harder final period.
\item \textbf{A decomposition of the oversampling leak.} We separate the artifact of scoring a rebalanced test set from genuine information flow between test and training data. The latter is present for every non-linear model, very large for some, and not detectable for logistic regression.
\item \textbf{Metrics tied to an alert budget.} We report the share of fraud cases and of fraud value caught when a fixed share of transactions can be reviewed, and show where in the ranking the inflation occurs.
\item \textbf{An explanation and a checklist.} We trace the inflation to fraud episodes that span the split, with transactions from the same card, minutes apart, on both sides. We condense the findings into a checklist for authors and reviewers and release code that reproduces every number in the paper.
\end{itemize}
